% chapters/ch11.tex
% 第~
% \ref{chap:11}~章 杂记与补充
% 源: 原书 PDF 第 352--373 页 (印刷页 349--370)

\chapter{杂记与补充}
\label{chap:11}

为了使全书的叙述更加紧凑, 本章收集了未纳入正文的若干补充内容.它们与第 1--7
章相联系.其中一些证明非常简略.有几个证明被略去, 有兴趣的读者可参阅参考文献.

\section{有限维与有限余维空间}
\label{sec:11.1}

众所周知, 任何 $p$ 维的有限维空间 $X$ 都同构于 $\R^p$.特别地, $X$ 完备,
$X$ 上的所有范数都等价, 且闭单位球 $B_X$ 是紧的.

\begin{proposition}
\label{prop:11.1}
设 $E$ 为 Banach 空间, $X\subset E$ 为有限维空间.则 $X$ 是闭的.
\end{proposition}

\begin{proof}
设 $(x_n)$ 为 $X$ 中的序列, 使得 $x_n\to x$ 在 $E$ 中.则 $(x_n)$ 是 $X$
中的 Cauchy 序列, 从而 $(x_n)$ 收敛于 $X$ 中的极限.因此 $x\in X$.
\end{proof}

\begin{proposition}
\label{prop:11.2}
假设 $X$ 有限维, $F$ 为 Banach 空间.则任何线性算子 $T:X\to F$ 必有界.
\end{proposition}

\begin{proof}
设 $(e_i)$ 为 $X$ 的基, 并把 $x$ 写成 $x=\sum_{i=1}^px_ie_i$.则
$Tx=\sum_{i=1}^px_iTe_i$, 从而
$\norm{Tx}\le\sum_{i=1}^p\abs{x_i}\norm{Te_i}\le(\max_i\norm{Te_i})\sum_{i=1}^p\abs{x_i}
\le C\norm{x}$.
\end{proof}

特别地, $X$ 上的所有线性泛函都连续.有限维空间 $X$ 的对偶空间 $X^\star$
也是有限维的, 且 $\dim X^\star=\dim X$.更精确地说, 若 $(e_i)$ 是 $X$ 的
基, 则把 $x\in X$ 写为 $x=\sum_{i=1}^px_ie_i$, 并令 $f_i(x)=x_i$,
$i=1,2,\dots,p$.显然泛函 $(f_i)$ 在 $X^\star$ 中线性无关且它们生成
$X^\star$.因此它们构成 $X^\star$ 的基.较不显然的是下述结论:

\begin{proposition}
\label{prop:11.3}
假设 $X$ 为 Banach 空间 (其中 $\dim X\le\infty$), 使得 $X^\star$ 有限维.
则 $X$ 有限维且 $\dim X=\dim X^\star$.
\end{proposition}

\begin{proof}
我们需要 Hahn--Banach 定理, 更精确地说是推论  \ref{cor:1.4}.设
$J:X\to X^{\star\star}$ 为 \S \ref{sec:1.3} 中定义的典范单射.由于
$\dim X^\star<\infty$, 由上面的讨论推出 $\dim X^{\star\star}<\infty$.但
$X$ 同构于 $J(X)\subset X^{\star\star}$, 从而
$\dim X\le\dim X^{\star\star}=\dim X^\star$.因此 $\dim X<\infty$, 再次
由上面的讨论推出 $\dim X^\star=\dim X$.
\end{proof}

\begin{proposition}
\label{prop:11.4}
设 $E$ 为 Banach 空间, $M\subset E$ 为闭子空间.假设 $X\subset E$ 为有限维
子空间.则 $(M+X)$ 是闭的.此外, $(M+X)$ 在 $E$ 中有补空间当且仅当 $M$ 有
补空间.
\end{proposition}

[警告: 回忆一般说来, 两个闭子空间的和不必是闭的; 参见例如习题 1.14.]

\begin{proof}
先额外假设 $M\cap X=\{0\}$.把 $u_n$ 写成 $u_n=x_n+y_n$, 其中 $x_n\in X$,
$y_n\in M$, 且 $u_n\to u$ 在 $E$ 中.我们断言 $(x_n)$ 有界.若不然, 则对
某个子序列 $n_k\to\infty$, 有 $\norm{x_{n_k}}\to\infty$.抽取进一步的
子序列, 可以假设在 $X$ 中 $\dfrac{x_{n_k}}{\norm{x_{n_k}}}\to\xi$, 满足
$\norm{\xi}=1$ (这里我们用到 $\dim X<\infty$ 的事实).从而
$\dfrac{y_{n_k}}{\norm{x_{n_k}}}=\dfrac{u_{n_k}}{\norm{x_{n_k}}}
-\dfrac{x_{n_k}}{\norm{x_{n_k}}}\to-\xi$; 此外, $\xi\in M$ (因为 $M$ 是
闭的).因此 $\xi\in M\cap X$, 我们必有 $\xi=0$.矛盾.于是我们证明了
$(x_n)$ 有界.抽取子序列, 可以假设 $x_{n_k}\to x$ 在 $X$ 中.则
$y_n\to u-x\in M$ (因为 $M$ 是闭的).因此 $u\in(M+X)$, 这完成了
$M\cap X=\{0\}$ 时 $(M+X)$ 闭的证明.

在一般情形下, 设 $\widetilde X$ 为 $(M\cap X)$ 在 $X$ 中的补空间 (这是
有限维的事情).显然 $\widetilde X$ 有限维, $M\cap\widetilde X=\{0\}$, 且
$M+\widetilde X=M+X$.我们已经证明了 $(M+\widetilde X)$ 是闭的, 从而
$(M+X)$ 也是闭的.

现在设 $M$ 在 $E$ 中有补空间, 记作 $N$.设 $P_M$ 与 $P_N$ 分别为到 $M$
与 $N$ 上的投影算子.由于 $P_N(X)$ 维数有限, 它在 $N$ 中有补空间, 记作
$\widetilde N$ (见 \S \ref{sec:2.4}).我们断言 $\widetilde N$ 是 $(M+X)$
在 $E$ 中的补空间.

首先我们有
\begin{equation*}
(M+X)\cap\widetilde N=\{0\}.
\end{equation*}
事实上, 若 $\tilde n=m+x$, 其中 $\tilde n\in\widetilde N$, $m\in M$,
$x\in X$, 则
\begin{equation*}
\tilde n=P_N\tilde n=P_N(m+x)=P_Nx\in P_N(X),
\end{equation*}
从而 $\tilde n\in\widetilde N\cap P_N(X)=\{0\}$.

其次, 我们有
\begin{equation*}
(M+X)+\widetilde N=E.
\end{equation*}
事实上, 任何 $\xi\in E$ 可以写成
\begin{equation*}
\xi=P_M\xi+P_N\xi,
\end{equation*}
而 $P_N\xi$ 可以进一步分解为
\begin{equation*}
P_N\xi=P_Nx+\tilde n,
\end{equation*}
其中 $x\in X$, $\tilde n\in\widetilde N$.但 $x=P_Mx+P_Nx$, 从而
\begin{equation*}
P_N\xi=(x-P_Mx)+\tilde n,
\end{equation*}
因此
\begin{equation*}
\xi=P_M\xi+x-P_Mx+\tilde n\in(M+X)+\widetilde N.
\end{equation*}

反之, 假设 $(M+X)$ 在 $E$ 中有补空间, 记作 $W$.如上, 设 $\widetilde X$
为 $(M\cap X)$ 在 $X$ 中的补空间.我们断言 $(W+\widetilde X)$ 是 $M$ 的
补空间.

首先我们有
\begin{equation*}
(W+\widetilde X)\cap M=\{0\}.
\end{equation*}
事实上, 若 $m\in M$ 可以写成 $m=w+\tilde x$, 其中 $w\in W$,
$\tilde x\in\widetilde X$, 则 $w=m-\tilde x$, 从而
$w\in(M+X)\cap W=\{0\}$.因此
$m=\tilde x\in(M\cap X)\cap\widetilde X=\{0\}$.

最后, 我们验证
\begin{equation*}
(W+\widetilde X)+M=E.
\end{equation*}
事实上, 只须验证
\begin{equation*}
(M+\widetilde X)=(M+X)
\end{equation*}
(因为 $W+(M+X)=E$).显然 $M+\widetilde X\subset M+X$ (因为
$\widetilde X\subset X$).反之, 任何 $x\in X$ 可以写成 $x=x_1+\tilde x$,
其中 $x_1\in M\cap X$, $\tilde x\in\widetilde X$.因此
$M+X\subset M+\widetilde X$.
\end{proof}

设 $M$ 为 Banach 空间 $E$ 的子空间.回忆: 若存在有限维空间 $X\subset E$
使得 $M+X=E$, 则称 $M$ 具有\emph{有限余维}.我们总可以假设 $M\cap X=\{0\}$
(否则在 $X$ 中选取 $M\cap X$ 的补空间).$M$ 的余维数按定义是这样的 $X$ 的
维数, 记作 $\operatorname{codim}M$ (它与 $X$ 的具体选取无关); 它与
$\dim(E/M)$ 一致.

[警告: 有限余维的子空间不必是闭的.例如, 若 $\dim E=\infty$, 取 $E$ 上的
任何不连续的线性泛函 $f$ (见习题 1.5).则 $M=f^{-1}(\{0\})$ 余维数为 1,
但 $M$ 不是闭的 (由命题  \ref{prop:1.5}); 事实上, $M$ 在 $E$ 中稠密.]

\begin{proposition}
\label{prop:11.5}
设 $E$ 为 Banach 空间, $M$ 为 $E$ 的有限余维的闭子空间.则包含 $M$ 的
任何子空间 $\widetilde M\subset E$ 必为闭的.
\end{proposition}

\begin{proof}
空间 $M$ 在 $\widetilde M$ 中有代数补空间, 记作 $X$.显然 $\dim X<\infty$,
且 $\widetilde M=X+M$.应用命题  \ref{prop:11.4}, 可知 $\widetilde M$ 是闭的.
\end{proof}

\begin{proposition}
\label{prop:11.6}
设 $E$ 为 Banach 空间, $M$ 为 $E$ 的有限余维的闭子空间.设 $D$ 为 $E$ 的
稠密子空间.则存在 $M$ 的补空间 $X$ 满足 $X\subset D$.
\end{proposition}

\begin{proof}
设 $d$ 为 $M$ 在 $E$ 中的余维数.若 $d=0$, 则 $M=E$, 可以取 $X=\{0\}$.
因此可以假设 $d\ge1$.固定任何满足 $x_1\notin M$ 的 $x_1\in D$; 这是可能
的, 否则 $D\subset M$ 蕴含 $E=\overline D\subset M\ne E$, 矛盾.令
$M_1=M+\R x_1$.则 $M_1$ 是闭的 (由命题  \ref{prop:11.4}), 且
$\operatorname{codim}M_1=d-1$.重复这一构造 $(d-1)$ 次, 即得到维数为 $d$
的子空间 $X\subset D$, 使得 $M+X=E$ 且 $M\cap X=\{0\}$.
\end{proof}

\begin{proposition}
\label{prop:11.7}
设 $E$ 为 Banach 空间, $G,L\subset E$ 为闭子空间.假设存在有限维空间
$X_1,X_2\subset E$ 使得
\begin{equation}
G+L+X_1=E,
\label{eq:11.1}
\end{equation}
\begin{equation}
G\cap L\subset X_2.
\label{eq:11.2}
\end{equation}
则 $G$ (相应地 $L$) 有补空间.
\end{proposition}

\begin{proof}
我们把证明分成两步.

\emph{第 1 步: 当 $X_2=\{0\}$ 时命题  \ref{prop:11.7} 的结论成立.}

设 $\widetilde X_1$ 为 $(G+L)\cap X_1$ 在 $X_1$ 中的补空间.由命题
 \ref{prop:11.4} 我们已知 $(L+\widetilde X_1)$ 是闭的.我们断言
$(L+\widetilde X_1)$ 是 $G$ 的补空间.

首先, 我们有
\begin{equation*}
G+(L+\widetilde X_1)=E.
\end{equation*}
事实上, 任何 $\xi\in E$ 可以写成 $\xi=g+\ell+h$, 其中 $g\in G$, $\ell\in L$,
$h\in X_1$, 而 $h$ 可以进一步分解为 $h=h_1+h_2$, 其中
$h_1\in(G+L)\cap X_1$, $h_2\in\widetilde X_1$.因此
$\xi\in G+L+\widetilde X_1$.

其次, 我们有
\begin{equation*}
G\cap(L+\widetilde X_1)=\{0\}.
\end{equation*}
事实上, 假设 $g=\ell+\tilde x_1$, 其中 $g\in G$, $\ell\in L$,
$\tilde x_1\in\widetilde X_1$.则 $\tilde x_1=g-\ell$, 从而
$\tilde x_1\in(G+L)\cap\widetilde X_1=\{0\}$.因此
$g=\ell\in G\cap L=\{0\}$ (这是第 1 步中的假设).

\emph{第 2 步: 一般情形.}

设 $\widetilde G$ 为 $(G\cap L)$ 在 $G$ 中的补空间, $\widetilde L$ 为
$(G\cap L)$ 在 $L$ 中的补空间 (注意, 由于 $G\cap L$ 有限维,
$\widetilde G$ 与 $\widetilde L$ 存在; 见 \S \ref{sec:2.4}).我们断言
\begin{equation}
(\widetilde G+\widetilde L)+(X_1+X_2)=E
\label{eq:11.3}
\end{equation}
以及
\begin{equation}
(\widetilde G\cap\widetilde L)=\{0\}.
\label{eq:11.4}
\end{equation}
这将完成本命题的证明.事实上, 由第 1 步推出 $\widetilde G$ 有补空间.
因此由命题  \ref{prop:11.4}, $G=\widetilde G+(G\cap L)$ 也有补空间.

\emph{\eqref{eq:11.3} 的验证.} 任何 $\xi\in E$ 可以写成
\begin{equation*}
\xi=g+\ell+x_1,\quad g\in G,\ \ell\in L,\ x_1\in X_1.
\end{equation*}
但 $g=\tilde g+h_1$, 其中 $\tilde g\in\widetilde G$, $h_1\in G\cap L$;
同理 $\ell=\tilde\ell+h_2$, 其中 $\tilde\ell\in\widetilde L$,
$h_2\in G\cap L$.因此
\begin{equation*}
\xi=(\tilde g+\tilde\ell)+x_1+(h_1+h_2)\in(\widetilde G+\widetilde L)+(X_1+X_2).
\end{equation*}

\emph{\eqref{eq:11.4} 的验证.} 设 $g\in\widetilde G\cap\widetilde L$.则
$g\in(G\cap L)\cap\widetilde L$ (因为 $\widetilde G\subset G$ 且
$\widetilde L\subset L$).但 $(G\cap L)\cap\widetilde L=\{0\}$.
\end{proof}
\section{商空间}
\label{sec:11.2}

设 $E$ 为 Banach 空间, $M$ 为闭子空间.我们在 $E$ 上考虑由 $x\sim y$
当且仅当 $x-y\in M$ 定义的等价关系.全体等价类的集合是一个向量空间, 记作
$E/M$, \index{shangkongjian@商空间 (quotient space)}
称为 $E$ (模 $M$) 的\emph{商空间}.把每个 $x\in E$ 映到其等价类
$[x]$ 的典范映射记作 $\pi:E\to E/M$.显然 $\pi$ 是满射线性算子.商空间
$E/M$ \index{shangfanshu@商范数 (quotient norm)}
配以\emph{商范数}
\begin{equation*}
\norm{[x]}_{E/M}=\norm{\pi(x)}_{E/M}
=\inf_{\substack{y\in E\\ y\in[x]}}\norm{y}
=\inf_{m\in M}\norm{x-m}.
\end{equation*}
容易看出 $\norm{[x]}_{E/M}$ 是 $E/M$ 上的范数 (为验证
$\norm{[x]}_{E/M}=0\Rightarrow[x]=0$, 须用到 $M$ 闭的事实).此外,
$\pi:E\to E/M$ 是有界算子, 且 $\norm{\pi}\le1$.若无混淆, 我们简写
$\norm{\cdot}$ 而不写 $\norm{\cdot}_{E/M}$.

\begin{proposition}
\label{prop:11.8}
配以范数 $\norm{\cdot}_{E/M}$ 的商空间 $E/M$ 是 Banach 空间.
\end{proposition}

\begin{proof}
设 $(\pi(x_k))$ 为 $E/M$ 中的 Cauchy 序列.我们须证明 $(\pi(x_k))$ 收敛;
由于 $(\pi(x_k))$ 是 Cauchy 序列, 只须证明某个子序列收敛.抽取子序列 (仍
记作 $(x_k)$), 可以假设对所有 $k$,
$\norm{\pi(x_{k+1})-\pi(x_k)}<\dfrac{1}{2^k}$ (见定理  \ref{thm:4.8} 的
证明).因此存在 $M$ 中的序列 $(m_k)$ 使得
$\norm{x_{k+1}-x_k-m_k}<\dfrac{1}{2^k}$.令 $m_k=\mu_{k+1}-\mu_k$, 其中
$\mu_1=0$, 且 $\mu_k\in M\ \forall k$.由于 $(x_k-\mu_k)$ 是 $E$ 中的
Cauchy 序列, 它在 $E$ 中收敛于极限 $\ell$.因此
$\pi(x_k)=\pi(x_k-\mu_k)$ 也在 $E/M$ 中收敛 (收敛于 $\pi(\ell)$).
\end{proof}

\begin{proposition}
\label{prop:11.9}
设 $M$ 为 $E$ 的闭子空间, 设 $\pi^\star:(E/M)^\star\to E^\star$ 为
$\pi:E\to E/M$ 的对偶算子.则 $R(\pi^\star)=M^\perp$, 更精确地说,
$\pi^\star$ 是由 $(E/M)^\star$ 到 $M^\perp$ 上的双射, 满足
\begin{equation*}
\norm{\pi^\star(\xi)}_{E^\star}=\norm{\xi}_{(E/M)^\star}
\quad\forall\xi\in(E/M)^\star.
\end{equation*}
特别地, $(E/M)^\star$ 与 $M^\perp$ 同构且等距.
\end{proposition}

\begin{proof}
设 $\xi\in(E/M)^\star$ 且 $x\in E$, 记
\begin{equation*}
\dual{\pi^\star(\xi)}{x}=\dual{\xi}{\pi(x)}.
\end{equation*}
若 $x\in M$, 则 $\pi(x)=0$, 从而 $\dual{\pi^\star(\xi)}{x}=0\ \forall x\in M$,
即 $\pi^\star(\xi)\in M^\perp$.

反之, 设 $f\in M^\perp$; 我们须证明对某个 $\xi\in(E/M)^\star$, 有
$f=\pi^\star(\xi)$.给定 $y\in E/M$, 把 $y$ 写成 $y=\pi(x)$, 其中 $x\in E$,
然后定义 $\xi(y)=\dual{f}{x}$.注意这个定义不依赖于 $x$ 的具体选取, 因为
$f\in M^\perp$.显然 $\xi$ 关于 $y$ 线性, 且对一切 $m\in M$, 有
$\abs{\xi(y)}\le\norm{f}_{E^\star}\norm{x-m}$.对所有 $m\in M$ 取下确界,
给出
$\abs{\xi(y)}\le\norm{f}_{E^\star}\norm{\pi(x)}_{E/M}
=\norm{f}_{E^\star}\norm{y}_{E/M}$.因此 $\xi\in(E/M)^\star$, 且显然
$\dual{\pi^\star(\xi)}{x}=\dual{\xi}{\pi(x)}=\dual{f}{x}\ \forall x\in E$,
即 $\pi^\star(\xi)=f$.此外,
$\norm{\xi}_{(E/M)^\star}\le\norm{f}_{E^\star}=\norm{\pi^\star(\xi)}_{E^\star}$.

另一方面, 我们知道
\begin{equation*}
\norm{\pi^\star(\xi)}_{E^\star}\le\norm{\xi}_{(E/M)^\star},
\end{equation*}
因为 $\norm{\pi^\star}=\norm{\pi}\le1$.从而
\begin{equation*}
\norm{\pi^\star(\xi)}_{E^\star}=\norm{\xi}_{(E/M)^\star}
\quad\forall\xi\in(E/M)^\star.
\end{equation*}
\end{proof}

设 $F,G$ 为 Banach 空间, $T\in\LL(F,G)$.考虑 $F$ 的闭子空间 $N(T)$,
商空间 $F/N(T)$, 以及典范映射 $\pi:F\to F/N(T)$.算子 $T$ 可以分解为
$T=\widetilde T\circ\pi$, 其中 $\widetilde T:F/N(T)\to G$; 事实上, 给定
$y\in F/N(T)$, 把 $y$ 写成 $y=\pi(x)$, 其中 $x\in F$, 并令
$\widetilde Ty=Tx$.显然 $\widetilde T$ 良定 (不依赖于 $x$ 的选取), 并且
是由 $F/N(T)$ 到 $R(T)$ 上的双射; 此外, $\norm{\widetilde T}=\norm{T}$.

现在考虑这个框架的一个特殊情形.设 $M$ 为 Banach 空间 $E$ 的闭子空间.
设 $T:E^\star\to M^\star$ 定义为
\begin{equation*}
T(f)=f|_M\quad\forall f\in E^\star.
\end{equation*}
则 $N(T)=M^\perp$, 且 (由 Hahn--Banach 定理) $R(T)=M^\star$.把上述结论
应用于 $F=E^\star$ 与 $G=M^\star$, 我们得到双射算子
$\widetilde T:E^\star/M^\perp\to M^\star$, 使得 $\widetilde T\circ\pi=T$.

\begin{proposition}
\label{prop:11.10}
对任何 Banach 空间 $E$ 与 $E$ 的任何闭子空间 $M$, 算子 $\widetilde T$ 是由
$E^\star/M^\perp$ 到 $M^\star$ 上的双射等距.
\end{proposition}

\begin{proof}
只须证明 $\widetilde T$ 是等距.给定任何 $f\in E^\star$, 考虑 $M$ 上的泛函
$f|_M$.由推论  \ref{cor:1.2} 我们知道, 存在泛函 $\tilde f\in E^\star$
使得 $\tilde f|_M=f|_M$ 且
$\norm{\tilde f}_{E^\star}=\norm{f|_M}_{M^\star}=\norm{T(f)}_{M^\star}$.
由于 $f-\tilde f\in M^\perp$, 我们有
\begin{equation*}
\norm{\pi(f)}_{E^\star/M^\perp}=\dist(f,M^\perp)
\le\norm{f-(f-\tilde f)}_{E^\star}=\norm{\tilde f}_{E^\star}
=\norm{Tf}_{M^\star}.
\end{equation*}
因此我们证明了
\begin{equation*}
\norm{\pi(f)}_{E^\star/M^\perp}\le\norm{Tf}_{M^\star}\quad\forall f\in E^\star.
\end{equation*}
但 $T=\widetilde T\circ\pi$, 从而
\begin{equation*}
\norm{\pi(f)}_{E^\star/M^\perp}\le\norm{\widetilde T(\pi(f))}_{M^\star}
\quad\forall f\in E^\star,
\end{equation*}
即
\begin{equation*}
\norm{y}_{E^\star/M^\perp}\le\norm{\widetilde T(y)}_{M^\star}
\quad\forall y\in E^\star/M^\perp.
\end{equation*}
另一方面, 显然有
\begin{equation*}
\norm{(\widetilde T\circ\pi)(f)}_{M^\star}=\norm{T(f)}_{M^\star}\le\norm{f}_{E^\star}
\quad\forall f\in E^\star.
\end{equation*}
把 $f$ 换成 $(f-g)$ (其中 $g\in M^\perp$) 并在 $g\in M^\perp$ 上取下确界,
得到
\begin{equation*}
\norm{\widetilde T(\pi(f))}_{M^\star}\le\norm{\pi(f)}_{E^\star/M^\perp}
\quad\forall f\in E^\star,
\end{equation*}
即
\begin{equation*}
\norm{\widetilde T(y)}_{M^\star}\le\norm{y}_{E^\star/M^\perp}
\quad\forall y\in E^\star/M^\perp.
\end{equation*}
由此得出 $\widetilde T$ 是等距.
\end{proof}

商空间 $E/M$ 继承了空间 $E$ 的许多性质, 例如自反性与一致凸性.

\begin{proposition}
\label{prop:11.11}
假设 $E$ 为自反 Banach 空间, $M$ 为闭子空间.则 $E/M$ 自反.
\end{proposition}

\begin{proof}
我们知道 $E^\star$ 自反 (见推论  \ref{cor:3.21}), 从而 $M^\perp$ 也自反
(作为 $E^\star$ 的闭子空间; 见命题  \ref{prop:3.20}).另一方面, $M^\perp$
同构于 $(E/M)^\star$ (由命题  \ref{prop:11.9}).因此 $(E/M)^\star$ 自反,
再次由推论  \ref{cor:3.21}, $E/M$ 也自反.
\end{proof}

\begin{proposition}
\label{prop:11.12}
假设 $E$ 为一致凸 Banach 空间, $M$ 为闭子空间.则 $E/M$ 一致凸.
\end{proposition}

\begin{proof}
设 $\pi(x),\pi(y)\in E/M$ 满足 $\norm{\pi(x)}\le1$, $\norm{\pi(y)}\le1$,
且 $\norm{\pi(x)-\pi(y)}>\varepsilon$.由于 $E$ 自反, 我们知道 (见推论
 \ref{cor:3.23}) 存在 $m_1\in M$ 与 $m_2\in M$ 使得
$\norm{x-m_1}\le1$ 且 $\norm{y-m_2}\le1$.此外, 对所有 $m\in M$,
$\norm{(x-y)-m}>\varepsilon$.由 $E$ 的一致凸性推出
\begin{equation*}
\norm{\frac{(x-m_1)+(y-m_2)}{2}}<1-\delta,
\end{equation*}
从而
\begin{equation*}
\norm{\frac{\pi(x)+\pi(y)}{2}}<1-\delta.
\end{equation*}
\end{proof}

\begin{proposition}
\label{prop:11.13}
设 $E$ 为 Banach 空间, $M\subset E$ 为闭子空间.则

(a) $\dim M<\infty$ 当且仅当 $\operatorname{codim}M^\perp<\infty$, 且在这时
$\dim M=\operatorname{codim}M^\perp$,

(b) $\operatorname{codim}M<\infty$ 当且仅当 $\dim M^\perp<\infty$, 且在这时
$\operatorname{codim}M=\dim M^\perp$.
\end{proposition}

\begin{proof}
(a) 由命题  \ref{prop:11.10} 我们知道, $E^\star/M^\perp$ 总是与 $M^\star$
同构.因此
$\dim M^\star<\infty\iff\dim(E^\star/M^\perp)<\infty$.由命题
 \ref{prop:11.3} 我们知道,
$\dim M<\infty\iff\dim M^\star<\infty$ 且此时 $\dim M=\dim M^\star$.
另一方面,
$\dim(E^\star/M^\perp)<\infty\iff\operatorname{codim}M^\perp<\infty$.
因此 $\dim M<\infty\iff\operatorname{codim}M^\perp<\infty$ 且
$\dim M=\dim M^\star=\dim(E^\star/M^\perp)=\operatorname{codim}M^\perp$.

(b) 命题  \ref{prop:11.9} 给出
$\dim M^\perp<\infty\iff\dim(E/M)^\star<\infty$.再次利用命题
 \ref{prop:11.3}, 这等价于 $\dim(E/M)<\infty$, 即
$\operatorname{codim}M<\infty$.此时
$\dim M^\perp=\dim(E/M)^\star=\dim(E/M)=\operatorname{codim}M$.
\end{proof}

一个``对偶''断言部分成立.

\begin{proposition}
\label{prop:11.14}
设 $N\subset E^\star$ 为闭子空间.则 $\dim N<\infty$ 当且仅当
$\operatorname{codim}N^\perp<\infty$, 且在这时
$\dim N=\operatorname{codim}N^\perp$.还成立
$\dim N^\perp\le\operatorname{codim}N$, 但可能出现
$\dim N^\perp<\operatorname{codim}N<\infty$ 的情形.
\end{proposition}

\begin{proof}
回忆
\begin{equation*}
N^\perp=\{x\in E;\ \dual{f}{x}=0\ \forall f\in N\}.
\end{equation*}
显然 $N\subset N^{\perp\perp}$; 但可能出现 $N\ne N^{\perp\perp}$ (见第 1
章注记 6).例如, 取 $\xi\in E^{\star\star}\setminus E$, 并令
$N=\xi^{-1}(\{0\})=\{f\in E^\star;\ \dual{\xi}{f}=0\}$.则 $N$ 是 $E^\star$
的余维数为 1 的闭子空间 (即 $N$ 是超平面).然而, $N^\perp=\{0\}$ (因为
由引理  \ref{lem:3.2}, $N$ 在 $E^{\star\star}$ 中的正交补为 $\R\xi$, 从而
$N$ 在 $E$ 中的正交补 $N^\perp$ 简化为 $\{0\}$).在这种情形下,
$N^\perp=\{0\}$, 从而 $N^{\perp\perp}=E^\star$, 且 $\dim N^\perp=0$,
而 $\operatorname{codim}N=1$.

现在我们回到一般情形.由于 $N\subset N^{\perp\perp}$, 我们有
$\operatorname{codim}N^{\perp\perp}\le\operatorname{codim}N\le\infty$.令
$M=N^\perp\subset E$.由命题  \ref{prop:11.10} 我们有
\begin{equation*}
\operatorname{codim}M^\perp=\dim(E^\star/M^\perp)=\dim M^\star,
\end{equation*}
从而 $\operatorname{codim}N^{\perp\perp}=\dim M\le\infty$.因此
\begin{equation*}
\dim N^\perp\le\operatorname{codim}N\le\infty.
\end{equation*}

现在证明 $\dim N<\infty\Rightarrow\operatorname{codim}N^\perp<\infty$ 且
$\operatorname{codim}N^\perp=\dim N$.

我们先断言 $N^{\perp\perp}=N$.我们已知 $N\subset N^{\perp\perp}$.设
$f_1,f_2,\dots,f_p$ 为 $N$ 的基, 设 $f\in N^{\perp\perp}$.由于在
$N^\perp=\{x\in E;\ \dual{f_i}{x}=0\ \forall i\}$ 上 $f=0$, 我们可以应用
引理  \ref{lem:3.2} 并得出 $f=\sum\lambda_if_i$.因此
$N^{\perp\perp}\subset N$.如上, 令 $M=N^\perp$.由于 $\dim M^\perp<\infty$,
由命题  \ref{prop:11.13} 我们推出 $\operatorname{codim}M<\infty$ 且
$\operatorname{codim}M=\dim M^\perp$, 即
$\operatorname{codim}N^\perp=\dim N$.

反之, 假设 $\operatorname{codim}N^\perp<\infty$, 再次令 $M=N^\perp$, 于是
$\operatorname{codim}M<\infty$.再次应用命题  \ref{prop:11.13} 给出
$\dim M^\perp<\infty$, 即 $\dim N^{\perp\perp}<\infty$.由于
$N\subset N^{\perp\perp}$, 我们推出 $\dim N<\infty$, 这回到了前面的情形.
因此 $\dim N=\operatorname{codim}N^\perp$.
\end{proof}
\section{若干经典的序列空间}
\label{sec:11.3}

给定序列 $x=(x_1,x_2,\dots,x_k,\dots)$, 令
\begin{align*}
\norm{x}_p&=\left(\sum_{k=1}^\infty\abs{x_k}^p\right)^{1/p},\quad 1\le p<\infty,\\
\norm{x}_\infty&=\sup_k\abs{x_k},
\end{align*}
并考虑相应的空间
\begin{align*}
\ell^p&=\{x;\ \norm{x}_p<\infty\},\quad 1\le p<\infty,\\
\ell^\infty&=\{x;\ \norm{x}_\infty<\infty\},
\end{align*}
它们关于 $\ell^p$ (相应地 $\ell^\infty$) 范数是 Banach 空间.这可以直接
证明 (且相当容易); 也可以应用定理  \ref{thm:4.8} 于配以计数测度
$\mu(E)=E\subset\N$ 中的点数的 $\Omega=\N$.下面提到的许多性质都是第~
\ref{chap:4}~章
中一般结果的推论.为了方便读者, 我们也给出一些直接的证明.

$\ell^\infty$ 有两个有趣的子空间:
\begin{align*}
c&=\left\{x;\ \lim_{k\to\infty}x_k\ \text{存在}\right\},\\
c_0&=\left\{x;\ \lim_{k\to\infty}x_k=0\right\}.
\end{align*}
它们都配以 $\ell^\infty$ 范数.显然 $c_0\subset c\subset\ell^\infty$, 其中
$c_0$ 在 $c$ 中闭, $c$ 在 $\ell^\infty$ 中闭.

Hölder 不等式采取如下形式:
\begin{equation}
\left|\sum_{k=1}^\infty x_ky_k\right|\le\norm{x}_p\norm{y}_{p'}
\quad\forall x\in\ell^p,\ \forall y\in\ell^{p'}\ \text{满足}\
\frac1p+\frac{1}{p'}=1.
\label{eq:11.5}
\end{equation}
空间 $\ell^2$ 是配以标量积
\begin{equation*}
(x,y)=\sum_{k=1}^\infty x_ky_k
\end{equation*}
的 Hilbert 空间.显然 $\ell^p\subset c_0$, 满足
\begin{equation*}
\norm{x}_\infty\le\norm{x}_p\quad\forall p,\ 1\le p<\infty,\ \forall x\in\ell^p,
\end{equation*}
这给出当 $1\le p\le q\le\infty$ 时 $\ell^p\subset\ell^q$, 满足
\begin{equation*}
\norm{x}_q\le\norm{x}_p\quad\forall x\in\ell^p.
\end{equation*}

\begin{proposition}
\label{prop:11.15}
对 $1<p<\infty$, 空间 $\ell^p$ 是自反的, 甚至是一致凸的.
\end{proposition}

\begin{proof}
应用定理  \ref{thm:4.10} 与习题 4.12 于 $\Omega=\N$.
\end{proof}

\begin{proposition}
\label{prop:11.16}
空间 $c$, $c_0$ 以及 $\ell^p$ ($1\le p<\infty$) 是可分的.
\end{proposition}

\begin{proof}
令
\begin{equation*}
D=\{x=(x_k);\ x_k\in\Q\ \forall k,\ \text{且对足够大的}\ k\ \text{有}\ x_k=0\}.
\end{equation*}
显然 $D$ 是可数的; 此外, 当 $1\le p<\infty$ 时 $D$ 在 $\ell^p$ 中稠密, 且
在 $c_0$ 中稠密.集合 $D+\lambda(1,1,1,\dots)$ ($\lambda\in\Q$) 是可数的,
且在 $c$ 中稠密.
\end{proof}

\begin{proposition}
\label{prop:11.17}
空间 $\ell^\infty$ 不可分.
\end{proposition}

\begin{proof}
假设 $A\subset\ell^\infty$ 是可数的.我们将验证 $A$ 不能在 $\ell^\infty$
中稠密.记 $A=(a_k)$, 其中每个 $a_k\in\ell^\infty$, 从而
$a_k=(a_k^1,a_k^2,\dots)$.对每个整数 $k$, 令
\begin{equation*}
b_k=\begin{cases}
a_k^k+1 & \text{若}\ \abs{a_k^k}\le1,\\
0 & \text{若}\ \abs{a_k^k}>1.
\end{cases}
\end{equation*}
注意 $b=(b_k)\in\ell^\infty$ 且 $\abs{b_k-a_k^k}\ge1\ \forall k$.因此,
\begin{equation*}
\norm{b-a_k}_\infty\ge\abs{b_k-a_k^k}\ge1\quad\forall k,
\end{equation*}
从而 $b\notin\overline A$.
\end{proof}

\begin{proposition}
\label{prop:11.18}
设 $1\le p<\infty$.给定任何 $\varphi\in(\ell^p)^\star$, 存在唯一的
$u\in\ell^{p'}$ 使得
\begin{equation*}
\dual{\varphi}{x}=\sum_{k=1}^\infty u_kx_k\quad\forall x\in\ell^p.
\end{equation*}
此外,
\begin{equation*}
\norm{u}_{p'}=\norm{\varphi}_{(\ell^p)^\star}.
\end{equation*}
\end{proposition}

\begin{proof}
令 $e_k=(0,0,\dots,\underset{(k)}{1},0,0,\dots)$.令 $u_k=\varphi(e_k)$.
我们断言 $u=(u_k)\in\ell^{p'}$ 且
\begin{equation}
\norm{u}_{p'}\le\norm{\varphi}_{(\ell^p)^\star}.
\label{eq:11.6}
\end{equation}
当 $p=1$ 时不等式 \eqref{eq:11.6} 是显然的, 因为
\begin{equation*}
\abs{u_k}\le\norm{\varphi}_{(\ell^1)^\star}\norm{e_k}_1
\le\norm{\varphi}_{(\ell^1)^\star}\quad\forall k.
\end{equation*}
现在转向 $1<p<\infty$ 的情形.固定整数 $N$.则对每个
$x=(x_1,x_2,\dots,x_N,0,0,\dots)$, 我们有
\begin{equation}
\sum_{k=1}^N u_kx_k=\varphi\left(\sum_{k=1}^N x_ke_k\right)
\le\norm{\varphi}_{(\ell^p)^\star}\norm{x}_p.
\label{eq:11.7}
\end{equation}
选取 $x_k=\abs{u_k}^{p'-2}u_k$ 给出
\begin{equation*}
\left(\sum_{k=1}^N\abs{u_k}^{p'}\right)^{1/p'}
\le\norm{\varphi}_{(\ell^p)^\star}.
\end{equation*}
令 $N\to\infty$, 可见 $u\in\ell^{p'}$ 且 \eqref{eq:11.6} 成立.此外,
\begin{equation*}
\varphi(x)=\sum_{k=1}^\infty u_kx_k\quad\forall x\in D,
\end{equation*}
其中 $D$ 是命题  \ref{prop:11.16} 的证明中定义的.由于 $D$ 在 $\ell^p$
中稠密, 我们得到
\begin{equation*}
\varphi(x)=\sum_{k=1}^\infty u_kx_k\quad\forall x\in\ell^p.
\end{equation*}
由 Hölder 不等式得出
\begin{equation*}
\abs{\varphi(x)}\le\norm{u}_{p'}\norm{x}_p\quad\forall x\in\ell^p,
\end{equation*}
从而 $\norm{\varphi}_{(\ell^p)^\star}\le\norm{u}_{p'}$.结合 \eqref{eq:11.6},
我们得到 $\norm{\varphi}_{(\ell^p)^\star}=\norm{u}_{p'}$.

$u$ 的唯一性是显然的.
\end{proof}

\begin{proposition}
\label{prop:11.19}
给定任何 $\varphi\in(c_0)^\star$, 存在唯一的 $u\in\ell^1$ 使得
\begin{equation*}
\dual{\varphi}{x}=\sum_{k=1}^\infty u_kx_k\quad\forall x\in c_0.
\end{equation*}
此外,
\begin{equation*}
\norm{u}_1=\norm{\varphi}_{(c_0)^\star}.
\end{equation*}
\end{proposition}

\begin{proof}
这是命题  \ref{prop:11.18} 的证明的一个容易的改编 (取 $p=\infty$ 且
$p'=1$); 证明的最后部分成立是因为 $D$ 在 $c_0$ 中稠密 (但在 $\ell^\infty$
中不稠密).
\end{proof}

\begin{proposition}
\label{prop:11.20}
给定 $\varphi\in(c)^\star$, 存在唯一的对 $(u,\lambda)\in\ell^1\times\R$
使得
\begin{equation*}
\dual{\varphi}{x}=\sum_{k=1}^\infty u_kx_k+\lambda\lim_{k\to\infty}x_k
\quad\forall x\in c.
\end{equation*}
此外,
\begin{equation*}
\norm{u}_1+\abs{\lambda}=\norm{\varphi}_{(c)^\star}.
\end{equation*}
\end{proposition}

\begin{proof}
把命题  \ref{prop:11.19} 应用于 $\varphi|_{c_0}$, 我们找到某个
$u\in\ell^1$ 使得
\begin{equation*}
\varphi(y)=\sum_{k=1}^\infty u_ky_k\quad\forall y\in c_0.
\end{equation*}
若 $x\in c$, 把 $x$ 写成 $x=y+ae$, 其中 $e=(1,1,1,\dots)$,
$a=\lim_{k\to\infty}x_k$, 且 $y\in c_0$.则
\begin{equation*}
\varphi(x)=\sum_{k=1}^\infty u_ky_k+a\varphi(e)
=\sum_{k=1}^\infty u_k(x_k-a)+a\varphi(e)
=\sum_{k=1}^\infty u_kx_k+\lambda a,
\end{equation*}
其中 $\lambda=\varphi(e)-\sum_{k=1}^\infty u_k$.

反之, 给定任何 $u\in\ell^1$ 与 $\lambda\in\R$, 泛函
\begin{equation}
\varphi(x)=\sum_{k=1}^\infty u_kx_k+\lambda\lim_{k\to\infty}x_k,
\quad x\in c,
\label{eq:11.8}
\end{equation}
定义了 $(c)^\star$ 的一个元素.我们断言
\begin{equation}
\norm{\varphi}_{(c)^\star}=\norm{u}_1+\abs{\lambda}.
\label{eq:11.9}
\end{equation}
显然
\begin{equation}
\norm{\varphi}_{(c)^\star}\le\norm{u}_1+\abs{\lambda}.
\label{eq:11.10}
\end{equation}
在 \eqref{eq:11.8} 中选取 $x=(x_k)$, 其中 $N$ 为固定的整数, 且
\begin{equation*}
x_k=\begin{cases}
\operatorname{sign}(u_k), & 1\le k\le N,\\
\operatorname{sign}(\lambda), & k>N,
\end{cases}
\end{equation*}
给出
\begin{equation*}
\varphi(x)=\sum_{k=1}^N\abs{u_k}+\operatorname{sign}(\lambda)
\sum_{k=N+1}^\infty u_k+\abs{\lambda}\le\norm{\varphi}_{(c)^\star}.
\end{equation*}
令 $N\to\infty$, 我们得到
\begin{equation*}
\norm{u}_1+\abs{\lambda}\le\norm{\varphi}_{(c)^\star},
\end{equation*}
结合 \eqref{eq:11.10}, 给出 \eqref{eq:11.9}.
\end{proof}

\begin{proposition}
\label{prop:11.21}
空间 $\ell^1$, $\ell^\infty$, $c$ 与 $c_0$ 不是自反的.
\end{proposition}

\begin{proof}
由命题  \ref{prop:11.19} 与  \ref{prop:11.18} 我们知道, $(c_0)^\star$ 是
$\ell^1$, 且 $(\ell^1)^\star$ 是 $\ell^\infty$.因此由 $c_0$ 到 $\ell^\infty$
中的恒等映射对应于 \S \ref{sec:1.3} 中定义的典范单射 $J:c_0\to(c_0)^{\star\star}$.
由于它不是满射, 我们得出 $c_0$ 不是自反的.应用推论  \ref{cor:3.21}, 我们
推出 $\ell^1$ 与 $\ell^\infty$ 也不是自反的.此外, $c$ 不可能是自反的;
否则, 作为 $c$ 的闭子空间的 $c_0$ 由命题  \ref{prop:3.20} 必自反.
\end{proof}

下表总结了上面讨论的主要性质:

\begin{table}[htbp]
	\centering
	\begin{tabular}{lccc}
		\hline
		空间 & 自反性 & 可分性 & 对偶空间 \\
		\hline
		$\ell^p$ ($1<p<\infty$) & 是 & 是 & $\ell^{p'}$ \\
		$\ell^1$ & 否 & 是 & $\ell^\infty$ \\
		$c_0$ & 否 & 是 & $\ell^1$ \\
		$c$ & 否 & 是 & $\ell^1\times\R$ \\
		$\ell^\infty$ & 否 & 否 & 严格大于 $\ell^1$ \\
		\hline
	\end{tabular}
\end{table}
\section{\texorpdfstring{$\CC$ 上的 Banach 空间: 哪些相似, 哪些不同?}{C 上的 Banach 空间: 哪些相似, 哪些不同?}}
\label{sec:11.4}

在本节中, 我们始终假设 $E$ 是 $\CC$ 上的向量空间.当然, 只须考虑把乘标量
$\lambda x$ 中的 $\lambda$ 限制在 $\R$ 中 ($x\in E$), 我们就可以与 $E$
联系一个 $\R$ 上的向量空间; 为了与 $E$ 区别, 相应的 $\R$ 上的向量空间常记作
$E_\R$.

线性子空间 $M\subset E$ 是满足对一切 $\lambda\in\CC$ 及一切 $x,y\in M$ 有
$\lambda x\in M$ 与 $x+y\in M$ 的子集 $M$.当然, $E$ 的线性子空间 $M$ 也是
$E_\R$ 的线性子空间.但逆命题不成立.例如, $\R^2$ 中包含 $0$ 的一条直线
$L$ 是 $\R^2$ 的线性子空间.然而, 若我们把 $\CC$ 与 $\R^2$ 等同起来, 则
直线 $L$ 不再是 $\CC$ 的线性子空间, 因为 $iL$ (即 $L$ 旋转 $\pi/2$) 不包含
在 $L$ 中.

$E$ 上的范数按定义是取值于 $[0,+\infty)$ 的函数 $\norm{\cdot}$, 使得
$\norm{x}=0\iff x=0$, $\norm{\lambda x}=\abs{\lambda}\norm{x}\ \forall\lambda\in\CC,
\ \forall x\in E$, 以及 $\norm{x+y}\le\norm{x}+\norm{y}$.显然
$\norm{\cdot}$ 也是 $E_\R$ 上的范数, 但逆命题不成立.

$E$ 上的线性泛函是映射 $f:E\to\CC$, 满足对一切 $\lambda\in\CC$ 与一切
$x,y\in E$, $f(\lambda x)=\lambda f(x)$ 且 $f(x+y)=f(x)+f(y)$.对偶空间
$E^\star$ 是 $E$ 上的所有连续线性泛函组成的空间; $E^\star$ 是 $\CC$ 上的
向量空间, 配以范数
\begin{equation*}
\norm{f}_{E^\star}=\sup_{\substack{x\in E\\\norm{x}\le1}}\abs{f(x)}.
\end{equation*}
复数 $f(x)$ 也记作 $\dual{f}{x}$, 显然对一切 $\lambda,\mu\in\CC$ 与一切
$x\in E$, 有 $\dual{\lambda f}{\mu x}=\lambda\mu\dual{f}{x}$.复对偶
$E^\star$ 与实对偶 $E^\star_\R$ 之间的对应由下述简单而富有启发性的结果给出.

\begin{proposition}
\label{prop:11.22}
映射
\begin{equation*}
I:f\in E^\star\mapsto\operatorname{Re}f\in E^\star_\R
\end{equation*}
是由 $E^\star$ 到 $E^\star_\R$ 上的双射等距.
\end{proposition}

\begin{proof}
显然,
\begin{equation*}
\abs{\operatorname{Re}\dual{f}{x}}\le\abs{\dual{f}{x}}\le\norm{f}_{E^\star}\norm{x},
\end{equation*}
从而
\begin{equation}
\norm{I(f)}_{E^\star_\R}\le\norm{f}_{E^\star}.
\label{eq:11.11}
\end{equation}
$I$ 是单射也是显然的, 因为 $\operatorname{Re}\dual{f}{x}=0\ \forall x\in E$
蕴含 $\operatorname{Re}\dual{f}{ix}=0\ \forall x\in E$, 即
$\operatorname{Im}\dual{f}{x}=0\ \forall x\in E$, 从而 $f=0$.接下来我们
断言 $I$ 是满射.事实上, 给定 $\varphi\in E^\star_\R$, 令
\begin{equation}
f(x)=\varphi(x)-i\varphi(ix)\quad\forall x\in E
\label{eq:11.12}
\end{equation}
[警告: $\varphi(ix)$ 不等于 $i\varphi(x)$, 因为 $\varphi(x)$ 与
$\varphi(ix)$ 都属于 $\R$].容易验证 $f\in E^\star$, 即对一切
$\lambda\in\CC$ 与一切 $x\in E$, $f(\lambda x)=\lambda f(x)$ (请读者验证!),
且 $I(f)=\operatorname{Re}f=\varphi$.由 \eqref{eq:11.11} 我们有
$\norm{\varphi}_{E^\star_\R}\le\norm{f}_{E^\star}$.由 \eqref{eq:11.12}
还清楚地看出
\begin{equation*}
\abs{f(x)}\le\left(\abs{\varphi(x)}^2+\abs{\varphi(ix)}^2\right)^{1/2}
\le\sqrt{2}\norm{\varphi}_{E^\star_\R}\norm{x}
\end{equation*}
(因为 $\norm{ix}=\norm{x}$).但我们可以做得更好.设 $f(x)\ne0$, 并令
$\lambda=\dfrac{f(x)}{\abs{f(x)}}\in\CC$.则
\begin{equation*}
\abs{f(x)}=\frac1\lambda f(x)=f\left(\frac{x}{\lambda}\right)
=\varphi\left(\frac{x}{\lambda}\right)-i\varphi\left(\frac{ix}{\lambda}\right).
\end{equation*}
由于 $\abs{f(x)}\in\R$, $\varphi\left(\dfrac{x}{\lambda}\right)\in\R$, 且
$\varphi\left(\dfrac{ix}{\lambda}\right)\in\R$, 可见
$\varphi\left(\dfrac{ix}{\lambda}\right)=0$, 从而
$\abs{f(x)}=\varphi\left(\dfrac{x}{\lambda}\right)$.因此
\begin{equation*}
\abs{f(x)}\le\norm{\varphi}_{E^\star_\R}\norm{\frac{x}{\lambda}}
=\frac1{\abs{\lambda}}\norm{\varphi}_{E^\star_\R}\norm{x}
=\norm{\varphi}_{E^\star_\R}\norm{x}.
\end{equation*}
从而 $\norm{f}_{E^\star}\le\norm{\varphi}_{E^\star_\R}
=\norm{I(f)}_{E^\star_\R}$.结合 \eqref{eq:11.11}, 我们得出 $I$ 是等距.
\end{proof}

命题  \ref{prop:11.22} 意味着: 当我们处理 $\CC$ 上的向量空间时, 第 1--5 章
中几乎没有任何改变, 只不过我们需要对 Hahn--Banach 定理稍加小心 (见下文).
重大差别出现在第~
\ref{chap:6}~章中我们处理特征值与谱的时候.这在有限维中就已经能看
出来: 任何复系数 $n\times n$ 矩阵 $M$ 在 $\CC$ 中必有特征值; 但即使 $M$
的元素属于 $\R$, 它在 $\R$ 中也可能没有特征值.我们现在逐章叙述需要作出的
修改.

\runin{第~
\ref{chap:1}~章.} 我们挑选几个例子来说明: 某些断言保持不变, 而另一些断言
需要微小的修改.

\begin{proposition}
\label{prop:11.23}
设 $G\subset E$ 为线性子空间.若 $g:G\to\CC$ 为连续线性泛函, 则存在
延拓 $g$ 的 $f\in E^\star$, 使得
\begin{equation*}
\norm{f}_{E^\star}=\norm{g}_{G^\star}.
\end{equation*}
\end{proposition}

\begin{proof}
令 $\psi=\operatorname{Re}g$, 于是 $\psi$ 是 $G^\star_\R$ 的元素, 且
$\norm{\psi}_{G^\star_\R}=\norm{g}_{G^\star}$.由推论  \ref{cor:1.2}, 存在
延拓 $\psi$ 的 $\phi\in E^\star_\R$, 使得
$\norm{\phi}_{E^\star_\R}=\norm{\psi}_{G^\star_\R}$.应用命题
 \ref{prop:11.22}, 可见存在 $f\in E^\star$ 使得
$\phi=\operatorname{Re}f$ 且
$\norm{f}_{E^\star}=\norm{\phi}_{E^\star_\R}=\norm{\psi}_{G^\star_\R}
=\norm{g}_{G^\star}$.此外, 在 $G$ 上有
$\phi=\operatorname{Re}f=\psi=\operatorname{Re}g$, 即对一切 $x\in G$,
$\operatorname{Re}f(x)=\operatorname{Re}g(x)$; 用 $ix$ 替代 $x$ 给出对一切
$x\in G$, $\operatorname{Im}f(x)=\operatorname{Im}g(x)$, 从而在 $G$ 上
$f=g$.
\end{proof}

接下来, 我们叙述 Hahn--Banach 定理的几何形式之一.$E$ 中的\emph{闭实超平面}
$H$ 是形如
\begin{equation*}
H=\{x\in E;\ \operatorname{Re}\dual{f}{x}=\alpha\}
=[\operatorname{Re}f=\alpha]
\end{equation*}
的集合, 其中 $f\in E^\star$, $f\ne0$, 且 $\alpha\in\R$.我们再次提醒读者:
若 $\alpha=0$, 则 $H$ 是 $E_\R$ 的线性子空间, 但它不是 $E$ 在 $\CC$ 上的
线性子空间; 例如, 在 $E=\CC$ 中, $H$ 是一条直线 (而直线不是 $E$ 的复线性
子空间).若对一切 $x\in A$, $\operatorname{Re}\dual{f}{x}\le\alpha$, 且对
一切 $x\in B$, $\operatorname{Re}\dual{f}{x}\ge\alpha$, 则称 $H$ \emph{分离}
$A,B\subset E$.

\begin{proposition}
\label{prop:11.24}
设 $A,B\subset E$ 为 $E$ 的两个非空凸子集, 满足 $A\cap B=\emptyset$.假设
其中之一是开集.则存在分离 $A$ 与 $B$ 的闭实超平面.
\end{proposition}

\begin{proof}
把定理  \ref{thm:1.6} 应用于 $E_\R$, 得到在通常意义下分离 $A$ 与 $B$ 的
超平面 $H=[\phi=\alpha]$ ($\phi\in E^\star_\R$).然后利用命题
 \ref{prop:11.23} 断言对某个 $f\in E^\star$, 有 $\phi=\operatorname{Re}f$.
\end{proof}

$E$ 的线性子空间 $M$ 的正交空间 $M^\perp$ 的定义不变:
\begin{equation*}
M^\perp=\{f\in E^\star;\ \dual{f}{x}=0\ \forall x\in M\},
\end{equation*}
显然我们有 $M^\perp=\{f\in E^\star;\ \operatorname{Re}\dual{f}{x}=0\
\forall x\in M\}$ (因为我们可以用 $ix$ 代替 $x$).容易看出 $M^{\perp\perp}=M$.

给定函数 $\varphi:E\to(-\infty,+\infty]$, 我们在 $E^\star$ 上定义其共轭
函数 $\varphi^\star$:
\begin{equation*}
\varphi^\star(f)=\sup_{x\in E}\{\operatorname{Re}\dual{f}{x}-\varphi(x)\}.
\end{equation*}
采用显而易见的记号, 我们有
\begin{equation*}
\varphi^\star(f)=\varphi^\star_\R(If)\quad\forall f\in E^\star.
\end{equation*}

\begin{proposition}
\label{prop:11.25}
假设 $\varphi:E\to(-\infty,+\infty]$ 是凸的, 下半连续的, 且
$\varphi\not\equiv+\infty$.则 $\varphi^{\star\star}=\varphi$.
\end{proposition}

\begin{proof}
有两种方法.或者把定理  \ref{thm:1.11} 应用于看作 $E_\R$ 上的
$\tilde\varphi=\varphi$, 并结合命题  \ref{prop:11.22}.或者重复定理
 \ref{thm:1.11} 的证明; 当使用 Hahn--Banach 定理时, 如上所述用实超平面
分离凸集 $A$ 与 $B$.
\end{proof}

示性函数 $I_K$ 的定义不变.若 $M$ 是 $E$ 的线性子空间且 $\varphi=I_M$, 则
\begin{equation*}
\varphi^\star(f)=\sup_{x\in M}\operatorname{Re}\dual{f}{x}=I_{M^\perp}.
\end{equation*}
事实上, 若 $f\in M^\perp$, 则对一切 $x\in M$ 有 $\dual{f}{x}=0$, 从而
$\varphi^\star(f)=0$.反之, 若 $f\notin M^\perp$, 则存在某个 $x_0\in M$
使得 $\dual{f}{x_0}\ne0$.必要时用 $ix_0$ 代替 $x_0$, 可以假设
$\operatorname{Re}\dual{f}{x_0}\ne0$.必要时用 $-x_0$ 代替 $x_0$, 可以
假设 $\operatorname{Re}\dual{f}{x_0}>0$, 从而
$\sup_{\lambda>0}\dual{f}{\lambda x_0}=+\infty$.

\runin{第~
\ref{chap:2}~章.} 所有断言均保持不变 (在推论~
\ref{cor:2.4} 与 \ref{cor:2.5} 中把 $\R$ 换成 $\CC$).
某些证明依赖于 $\R$ 的结构 (例如定理~
\ref{thm:2.16} 的证明中的公式 \eqref{eq:2.21}).它们容易
改编到 $\CC$; 另外, 也可以通过把 $\R$ 版本应用于 $E_\R$ 来建立 $\CC$ 上的
断言.

\runin{第~
\ref{chap:3}~章.} 所有断言均保持不变 (在引理~
\ref{lem:3.2} 与 \ref{lem:3.3} 中把 $\R$ 换成 $\CC$).
某些证明需要显而易见的修改 (例如命题~
\ref{prop:3.11} 的证明).

\runin{第~
\ref{chap:4}~章.} 完全不变.

\runin{第~
\ref{chap:5}~章.} $\CC$ 上的 Hilbert 空间是配以标量积 $(u,v)\in\CC$ 的 $\CC$
上的向量空间.这是由 $H\times H$ 到 $\CC$ 的映射, 满足
\begin{align*}
(u,v)&=\overline{(v,u)} &&\forall u,v\in H,\\
\text{对每个}\ v\in H,\ u&\mapsto(u,v)\ \text{是线性的},\\
(u,u)&>0 &&\forall u\ne0.
\end{align*}
特别地, 我们有
\begin{equation*}
(\lambda u,\mu v)=\lambda\bar\mu(u,v)\quad\forall\lambda,\mu\in\CC,\
\forall u,v\in H.
\end{equation*}
量 $\abs{u}=(u,u)^{1/2}$ 是一个范数; 我们有
\begin{equation*}
\abs{u+v}^2=\abs{u}^2+2\operatorname{Re}(u,v)+\abs{v}^2\quad\forall u,v\in H,
\end{equation*}
且 Cauchy--Schwarz 不等式采取如下形式:
\begin{equation*}
\abs{(u,v)}\le\abs{u}\abs{v}\quad\forall u,v\in H.
\end{equation*}
一个典型的例子是配以标量积
\begin{equation*}
(u,v)=\int_\Omega u(x)\overline{v(x)}d\mu
\end{equation*}
的 $L^2(\Omega;\CC)$.

$\R$ 上的 Hilbert 空间与 $\CC$ 上的 Hilbert 空间之间的联系如下: 假设 $H$
是 $\CC$ 上的 Hilbert 空间.则配以标量积 $\operatorname{Re}(u,v)$ 的 $H_\R$
成为 $\R$ 上的 Hilbert 空间.因此第~
\ref{chap:5}~章的所有断言对 $H_\R$ 都适用.下面是
一些例子.

\begin{proposition}
\label{prop:11.26}
设 $K\subset H$ 为非空闭凸集.则对每个 $f\in H$, 存在唯一的元素 $u\in K$
使得
\begin{equation*}
\abs{f-u}=\min_{v\in K}\abs{f-v}=\dist(f,K).
\end{equation*}
此外, $u$ 由下述性质表征:
\begin{equation*}
u\in K\quad\text{且}\quad\operatorname{Re}(f-u,v-u)\le0\quad\forall v\in K.
\end{equation*}
\end{proposition}

\begin{proposition}
\label{prop:11.27}
给定任何 $\phi\in H^\star$, 存在唯一的 $f\in H$ 使得
\begin{equation*}
\phi(u)=(u,f)\quad\forall u\in H.
\end{equation*}
此外,
\begin{equation*}
\abs{f}=\norm{\phi}_{H^\star}.
\end{equation*}
\end{proposition}

\begin{proof}
把定理  \ref{thm:5.5} 应用于 $H_\R$ 中的 $\operatorname{Re}\phi$, 我们找到
某个 $f\in H$ 使得
\begin{equation*}
\operatorname{Re}\phi(u)=\operatorname{Re}(u,f)\quad\forall u\in H.
\end{equation*}
把它应用于 $iu$ 给出 $\operatorname{Im}\phi(u)=\operatorname{Im}(u,f)$,
从而 $\phi(u)=(u,f)\ \forall u\in H$.
\end{proof}

现在考虑满足下述条件的函数 $a(u,v):H\times H\to\CC$:
\begin{equation}
\forall v\in H,\ u\mapsto a(u,v)\ \text{是线性的, 且}\
\forall u\in H,\ v\mapsto\overline{a(u,v)}\ \text{是线性的},
\label{eq:11.13}
\end{equation}
\begin{equation}
a\ \text{连续, 即}\ \abs{a(u,v)}\le C\abs{u}\abs{v}\quad\forall u,v\in H,
\label{eq:11.14}
\end{equation}
\begin{equation}
a\ \text{强制, 即}\ \operatorname{Re}a(u,u)\ge\alpha\abs{u}^2
\quad\forall u\in H,\ \text{对某个}\ \alpha>0.
\label{eq:11.15}
\end{equation}

\begin{proposition}
\label{prop:11.28}
假设 $a$ 满足 \eqref{eq:11.13}, \eqref{eq:11.14} 与 \eqref{eq:11.15}.
设 $K$ 为 $H$ 中的非空闭凸集.则给定任何 $\phi\in H^\star$, 存在唯一的
$u\in K$ 使得
\begin{equation}
\operatorname{Re}a(u,v-u)\ge\operatorname{Re}\dual{\phi}{v-u}
\quad\forall v\in K.
\label{eq:11.16}
\end{equation}
此外, 若对一切 $v,w\in H$ 有 $a(v,w)=\overline{a(w,v)}$, 则 $u$ 由下述
性质表征:
\begin{equation*}
u\in K\quad\text{且}\quad
\frac12a(u,u)-\operatorname{Re}\dual{\phi}{u}
=\min_{v\in K}\left\{\frac12a(v,v)-\operatorname{Re}\dual{\phi}{v}\right\}.
\end{equation*}
\end{proposition}

当 $K=H$ 时, \eqref{eq:11.16} 变为 $a(u,v)=\dual{\phi}{v}\ \forall v\in H$.
特别地, 我们推出任何满足
\begin{equation}
\operatorname{Re}(Tu,u)\ge\alpha\abs{u}^2\quad\forall u\in H,\
\text{对某个}\ \alpha>0
\label{eq:11.17}
\end{equation}
的算子 $T\in\LL(H)$ 是由 $H$ 到自身的双射.有一个看上去稍微更一般的变体
(然而见下面的注记 1).

\begin{proposition}[Lax--Milgram]
\label{prop:11.29}
假设 $T\in\LL(H)$ 满足
\begin{equation}
\abs{(Tu,u)}\ge\alpha\abs{u}^2\quad\forall u\in H,\ \text{对某个}\ \alpha>0.
\label{eq:11.18}
\end{equation}
则 $T$ 是双射.
\end{proposition}

\begin{proof}
见第~
\ref{chap:5}~章注记 8.
\end{proof}

\begin{remark}
\label{rem:11.1}
显然 \eqref{eq:11.17} 蕴含 \eqref{eq:11.18}.反之, 假设 \eqref{eq:11.18}
成立.则存在某个满足 $\abs{\xi}=1$ 的 $\xi\in\CC$, 使得
\begin{equation}
\operatorname{Re}(\xi Tu,u)\ge\alpha\abs{u}^2\quad\forall u\in H.
\label{eq:11.19}
\end{equation}
事实上, 数值域
\begin{equation*}
W(T)=\{(Tu,u);\ u\in H,\ \abs{u}=1\}
\end{equation*}
是凸集 (由下面的命题  \ref{prop:11.33}).此外, 由 \eqref{eq:11.18} 我们
知道 $0\notin W(T)$, 且事实上 $\dist(0,W(T))\ge\alpha$.设 $p$ 表示 $0$
在 $W(T)$ 上的投影 (在 $\CC\simeq\R^2$ 中).经过平面内的一个把 $p$ 转到
$x$ 轴上的点 $(0,\abs{p})$ 的旋转 (即乘以 $\xi\in\CC$, $\abs{\xi}=1$)
之后, 我们得出 \eqref{eq:11.19} 成立.
\end{remark}

\runin{第~
\ref{chap:6}~章.} \S\ref{sec:6.1} 与 \S\ref{sec:6.2} 完全不变.主要区别出现在 \S\ref{sec:6.3} 中.

设 $E$ 为 $\CC$ 上的 Banach 空间, $T\in\LL(E)$.预解集定义为
\begin{equation*}
\rho(T)=\{\lambda\in\CC;\ (T-\lambda I)\ \text{是由}\ E\ \text{到}\ E\
\text{上的双射}\}.
\end{equation*}
谱是 $\rho(T)$ 的余集, 即 $\sigma(T)=\CC\setminus\rho(T)$.若相应的特征
子空间 $N(T-\lambda I)\ne\{0\}$, 则称数 $\lambda\in\CC$ 为\emph{特征值},
全体特征值的集合记作 $\operatorname{EV}(T)$.显然
$\operatorname{EV}(T)\subset\sigma(T)$.可能出现
$\operatorname{EV}(T)=\emptyset$ 的情形 (例如右移位算子
$Tu=(0,u_1,u_2,\dots)$).然而, $\sigma(T)$ 绝非空集.

\begin{proposition}
\label{prop:11.30}
谱 $\sigma(T)$ 是非空紧集, 且
\begin{equation*}
\sigma(T)\subset\{\lambda\in\CC;\ \abs{\lambda}\le\norm{T}\}.
\end{equation*}
\end{proposition}

\begin{proof}
主要的新颖之处在于 $\sigma(T)$ 非空.证明依赖于 $\CC$ 上的解析函数理论
(更精确地说是 Liouville 定理), 我们这里不予给出.有兴趣的读者可查阅
A.~Taylor--D.~Lay \bref{Taylor}{1}, W.~Rudin \bref{Rudin}{2}, 或
A.~Knapp \bref{Knapp}{2}.
\end{proof}

估计 $\abs{\lambda}\le\norm{T}\ \forall\lambda\in\sigma(T)$ 通常不是精确的.
例如, 在 $\CC^2$ 中算子 $T(u_1,u_2)=(u_2,0)$ 满足 $\sigma(T)=\{0\}$ 且
$\norm{T}=1$.最优界是用\emph{谱半径}给出的.我们已知 (见习题 6.23): 对每个
算子 $T\in\LL(E)$,
\begin{equation*}
r(T)=\lim_{n\to\infty}\norm{T^n}^{1/n}\ \text{存在},
\end{equation*}
且显然 $r(T)\le\norm{T}$; $r(T)$ \index{pubanjing@谱半径 (spectral radius)}
称为\emph{谱半径}.

\begin{proposition}
\label{prop:11.31}
对每个 $T\in\LL(E)$, 我们有
\begin{equation*}
r(T)=\max\{\abs{\lambda};\ \lambda\in\sigma(T)\}.
\end{equation*}
\end{proposition}

关于证明, 我们再次引证 A.~Taylor--D.~Lay \bref{Taylor}{1},
W.~Rudin \bref{Rudin}{2}, 或 A.~Knapp \bref{Knapp}{2}.论证通过 $\CC$ 上的
幂级数理论在很大程度上依赖于 $E$ 是 $\CC$ 上的 Banach 空间这一事实.当 $E$
是 $\R$ 上的 Banach 空间时, 我们只能说
$\max\{\abs{\lambda};\ \lambda\in\sigma(T)\}\le r(T)$, 且即使
$\sigma(T)$ 非空, 不等式也可以是严格的 (见习题 6.23).

实空间与复空间之间的另一个有趣区别与\index{puyingshedingli@谱映射定理 (spectral mapping theorem)}
所谓的\emph{谱映射定理}有关.先考虑
实情形: 设 $Q(t)=\sum_{k=0}^pa_kt^k$ 为实系数多项式, 设 $T\in\LL(E)$,
其中 $E$ 是 $\R$ 上的 Banach 空间.我们知道 (见习题 6.22)
\begin{equation}
Q(\operatorname{EV}(T))\subset\operatorname{EV}(Q(T))
\quad\text{且}\quad
Q(\sigma(T))\subset\sigma(Q(T)),
\label{eq:11.20}
\end{equation}
且这些包含关系可以是严格的 (但在 Hilbert 空间中当 $T^\star=T$ 时除外).
在复情形, 这些包含关系成为等式: 假设 $Q(t)$ 是复系数多项式, 设
$T\in\LL(E)$, 其中 $E$ 是 $\CC$ 上的 Banach 空间.

\begin{proposition}
\label{prop:11.32}
我们有
\begin{equation}
Q(\operatorname{EV}(T))=\operatorname{EV}(Q(T))
\label{eq:11.21}
\end{equation}
以及
\begin{equation}
Q(\sigma(T))=\sigma(Q(T)).
\label{eq:11.22}
\end{equation}
\end{proposition}

\begin{proof}
我们已知 \eqref{eq:11.20} 成立 (论证与习题 6.22 相同).反设包含关系是
严格的.则存在 $\mu\in\operatorname{EV}(Q(T))$ 使得
$\mu\notin Q(\operatorname{EV}(T))$.写出
\begin{equation*}
Q(t)-\mu=\alpha(t-t_1)(t-t_2)\dots(t-t_p),
\end{equation*}
其中 $\alpha\ne0$, 且 $t_i\notin\operatorname{EV}(T)\ \forall i$.此外,
存在某个 $x\ne0$ 使得 $Q(T)x=\mu x$.由于 $(T-t_1I)$ 是单射, 我们推出
$(T-t_2I)\dots(T-t_pI)x=0$, 重复同样的论证给出 $x=0$.矛盾.

类似地, 设 $\mu\in\sigma(Q(T))$ 满足 $\mu\notin Q(\sigma(T))$.如上写出
$Q(t)-\mu$, 其中 $t_i\notin\sigma(T)\ \forall i$.则 $Q(T)-\mu I$ 可以
写成双射算子的乘积.因此 $Q(T)-\mu I$ 是双射, 即
$\mu\in\rho(Q(T))$.矛盾.
\end{proof}

在 Hilbert 空间中, 研究谱的一个有用工具是数值域.设 $H$ 为 $\CC$ 上的
Hilbert 空间; \index{shuzhiyu@数值域 (numerical range)}
算子 $T\in\LL(H)$ 的\emph{数值域}定义为
\begin{equation*}
W(T)=\{(Tu,u);\ u\in H\ \text{且}\ \abs{u}=1\}.
\end{equation*}

\begin{proposition}
\label{prop:11.33}
我们有
\begin{equation*}
\sigma(T)\subset\overline{W(T)},
\end{equation*}
更精确地说, 若 $\lambda\notin\overline{W(T)}$, 则 $\lambda\in\rho(T)$,
满足
\begin{equation}
\norm{(T-\lambda I)^{-1}}\le\frac{1}{\dist(\lambda,W(T))}.
\label{eq:11.23}
\end{equation}
此外, $W(T)$ 是凸的.
\end{proposition}

\begin{proof}
假设 $\lambda\notin\overline{W(T)}$, 并令 $\alpha=\dist(\lambda,W(T))$.
对满足 $\abs{u}=1$ 的一切 $u\in H$, 有
\begin{equation*}
\abs{(Tu,u)-\lambda}\ge\alpha.
\end{equation*}
从而
\begin{equation*}
\abs{(Tu-\lambda u,u)}\ge\alpha\abs{u}^2\quad\forall u\in H.
\end{equation*}
应用 Lax--Milgram 定理 (命题  \ref{prop:11.29}), 我们得出 $(T-\lambda I)$
是双射, 且 $\abs{Tu-\lambda u}\ge\alpha\abs{u}\ \forall u\in H$, 即
$\norm{(T-\lambda I)^{-1}}\le1/\alpha$.

$W(T)$ 的凸性是由于 Toeplitz 与 Hausdorff 的一个违反直觉的事实.证明
参见 P.~R.~Halmos \bref{Halmos}{2}.
\end{proof}

一般说来, 数值域 $W(T)$ 可能比谱大得多.例如, 取 $H=\CC^2$ 与
$T(u_1,u_2)=(u_2,0)$, 我们有
$\operatorname{EV}(T)=\sigma(T)=\{0\}$, 而
$W(T)=\{\lambda\in\CC;\ \abs{\lambda}\le1/2\}$.然而, 若 $T$ 自伴, 或更
一般地为\index{zhengguisuanzi@正规算子 (normal operator)}
正规算子 (见下文), 则 $W(T)=\operatorname{conv}\sigma(T)$, 即
$\sigma(T)$ 的凸包 (见 P.~R.~Halmos \bref{Halmos}{2} 与下面的注记 2).

当 $H$ 是 $\CC$ 上的 Hilbert 空间且 $T\in\LL(H)$ 时, 对伴随算子 $T^\star$
的概念需要提醒读者注意.依照一般的步骤, 算子 $T\in\LL(H)$ 的伴随算子通过
关系式
\begin{equation*}
\dual{T^\star f}{u}_{H^\star,H}=\dual{f}{Tu}_{H^\star,H}
\quad\forall f\in H^\star,\ \forall u\in H
\end{equation*}
定义, 此时 $T^\star\in\LL(H^\star)$ (我们强调: 对一切 $\lambda\in\CC$ 及
一切 $f\in H^\star$, $T^\star(\lambda f)=\lambda T^\star f$).此外,
对一切 $\lambda\in\CC$ 有 $(\lambda T)^\star=\lambda T^\star$ (因为
$f:H\to\CC$ 是线性的).

另一方面, 我们也可以把 $H^\star$ 与 $H$ 等同起来 (通过命题
 \ref{prop:11.27} 中的同构), 并把 $T^\star$ 看作由 $H$ 到自身的算子, 它
通过关系式
\begin{equation*}
(Tu,v)=(u,T^\star v)\quad\forall u,v\in H
\end{equation*}
定义, 此时 $T^\star\in\LL(H)$ (我们强调: 对一切 $\lambda\in\CC$ 及一切
$v\in H$, $T^\star(\lambda v)=\lambda T^\star v$).然而, 我们现在有
\begin{equation}
(\lambda T)^\star=\bar\lambda T^\star\quad\forall\lambda\in\CC
\label{eq:11.24}
\end{equation}
(这容易验证).通常采用这一约定, 使得 $T^\star$ 与 $T$ 生活在同一个世界中:
可以比较 $T^\star$ 与 $T$, 作复合 $T^\star\circ T$ 等.

若 $T^\star=T$, 即
\begin{equation*}
(Tu,v)=(u,Tv)\quad\forall u,v\in H,
\end{equation*}
则称算子 $T\in\LL(H)$ 是\emph{自伴的} (或 \emph{Hermite 的}).若 $T$
自伴, 则对一切 $u\in H$ 有 $(Tu,u)=(u,Tu)=\overline{(Tu,u)}$, 从而对一切
$u\in H$ 有 $(Tu,u)\in\R$.特别地, 数值域 $W(T)$ 是 $\R$ 的子集, 从而
$\sigma(T)\subset\R$.

紧自伴算子的谱分解与第~
\ref{chap:6}~章中完全相同.

\begin{proposition}
\label{prop:11.34}
设 $H$ 为 $\CC$ 上的可分 Hilbert 空间, 设 $T$ 为紧自伴算子.则存在由 $T$
的特征向量组成的 Hilbert 基 (且相应的特征值为实数).
\end{proposition}

若算子 $T\in\LL(H)$ 满足 $T^\star\circ T=T\circ T^\star$, 则称它为
\emph{正规的}.当底层空间 $H$ 是 $\R$ 上的 Hilbert 空间时, 正规算子的各种
性质在问题 43 中作了讨论; 当 $H$ 是 $\CC$ 上的 Hilbert 空间时它们仍然成立.
但现在我们有更多结果:

\begin{proposition}
\label{prop:11.35}
设 $H$ 为 $\CC$ 上的 Hilbert 空间, 设 $T$ 为正规算子.则
\begin{equation}
\max\{\abs{\lambda};\ \lambda\in\sigma(T)\}=\norm{T}.
\label{eq:11.25}
\end{equation}
\end{proposition}

\begin{proof}
由于 $T$ 正规, 对每个整数 $p\ge1$ 有
\begin{equation*}
\norm{T^p}=\norm{T}^p.
\end{equation*}
这在 $H$ 是 $\R$ 上的 Hilbert 空间时已在问题 43 中证明, 同样的论证当 $H$
是 $\CC$ 上的 Hilbert 空间时仍然有效 (或者也可以把实结果应用于 $H_\R$ 上的
$T$).因此 $r(T)=\lim_{n\to\infty}\norm{T^n}^{1/n}=\norm{T}$.结合命题
 \ref{prop:11.31}, 给出 \eqref{eq:11.25}.
\end{proof}

\begin{proposition}
\label{prop:11.36}
设 $H$ 为 $\CC$ 上的可分 Hilbert 空间, 设 $T$ 为紧正规算子.则存在由 $T$
的特征向量组成的 Hilbert 基 (但相应的特征值不必为实数).
\end{proposition}

\begin{proof}
若 $T$ 正规, 则对任何 $\lambda\in\CC$, $(T-\lambda I)$ 也是正规的.因此
(如同问题 43 中那样) 我们有
$N(T-\lambda I)=N((T-\lambda I)^\star)=N(T^\star-\bar\lambda I)$.由此
推出当 $\lambda\ne\mu$ 时, $N(T-\lambda I)$ 与 $N(T-\mu I)$ 正交.我们
就可以完全像定理  \ref{thm:6.11} 的证明那样进行.我们在 $F^\perp$ 上得到
紧正规算子 $T_0$, 满足 $\sigma(T_0)=\{0\}$.我们不再援引推论  \ref{cor:6.10}
来得出 $T_0=0$, 而是应用命题  \ref{prop:11.35} 推出 $T_0=0$.正是在这里,
我们利用了 $H$ 是 $\CC$ 上的空间这一事实 (同样的结论在实空间中不成立).
\end{proof}

\begin{remark}
\label{rem:11.2}
当 $T$ 是紧正规算子时, 容易由命题  \ref{prop:11.36} 推出
$W(T)=\operatorname{conv}\sigma(T)$.事实上, 如命题  \ref{prop:11.36} 中
那样选取基 $(e_i)$.给定满足 $\abs{u}=1$ 的 $u\in H$, 把 $u$ 写成
$u=\sum u_ie_i$, 其中 $\sum\abs{u_i}^2=1$.则 $Tu=\sum\lambda_iu_ie_i$ 且
$(Tu,u)=\sum\lambda_i\abs{u_i}^2$.对任何正规算子 $T$ (不必紧),
$W(T)=\operatorname{conv}\sigma(T)$ 仍然成立; 参见 P.~R.~Halmos \bref{Halmos}{2}.
\end{remark}

设 $H$ 为 $\CC$ 上的 Hilbert 空间.若 $\abs{Tu}=\abs{u}\ \forall u\in H$,
则称算子 $T\in\LL(H)$ 为\emph{等距}; 若 $T$ 是等距且又是满射, 则称 $T$
为\emph{\index{yousuanzi@酉算子 (unitary operator)}
酉算子}.当底层空间 $H$ 是 $\R$ 上的 Hilbert 空间时, 等距与酉算子
的各种性质在问题 44 中作了讨论; 当 $H$ 是 $\CC$ 上的 Hilbert 空间时它们中
的大多数仍然成立, 但关于谱的断言除外, 它需要作如下修改:

\begin{proposition}
\label{prop:11.37}
设 $T$ 为等距.则
\begin{equation*}
\operatorname{EV}(T)\subset S^1=\{\lambda\in\CC;\ \abs{\lambda}=1\}.
\end{equation*}
若 $T$ 是酉算子, 则
\begin{equation*}
\sigma(T)\subset S^1,
\end{equation*}
若 $T$ 不是酉算子, 则
\begin{equation*}
\sigma(T)=\{\lambda\in\CC;\ \abs{\lambda}\le1\}.
\end{equation*}
\end{proposition}

证明是问题 44 第 6 问的解答中给出的证明的一个容易的改编.

若 $T^\star=-T$, 则\index{xiezibansuanzi@斜自伴算子 (skew-adjoint operator)}
称算子 $T\in\LL(H)$ 为\emph{斜自伴的} (或
\emph{反对称的}).显然, $T$ 斜自伴当且仅当 $iT$ 自伴 (这由 \eqref{eq:11.24}
推出).因此, 对任何斜自伴算子, 我们有
$\operatorname{EV}(T)\subset\sigma(T)\subset\overline{W(T)}\subset i\R$.
